\documentclass[journal]{IEEEtran}
\usepackage{amsmath,amssymb}
\usepackage{graphicx}
\usepackage{tikz}
\usetikzlibrary{shapes,arrows,positioning}
\usepackage{hyperref}
\usepackage{booktabs}

% Auto-generated by tools/gen_methodology.py — DO NOT EDIT MANUALLY
% Generated: 2026-03-05 07:57 UTC

\begin{document}

%% =========================================================================
\section{Methodology}
\label{sec:methodology}
%% =========================================================================

This section describes the proposed Integrity-Aware Planning (IAP) framework.
The system integrates LiDAR-inertial state estimation with GNSS tightly-coupled
fusion, online integrity monitoring, and receding-horizon planning whose cost
function penalises Protection Level (PL) violations.

%% -------------------------------------------------------------------------
\subsection{System Overview}
\label{sec:overview}
%% -------------------------------------------------------------------------

Figure~\ref{fig:system_flow} shows the overall data flow.

\begin{figure}[ht]
  \centering

  \includegraphics[width=\columnwidth]{docs/figures/system_flow.pdf}
  \caption{IAP system flowchart. Grey arrows show data flow; coloured blocks
           correspond to the modules described in the subsections below.}
  \label{fig:system_flow}
\end{figure}


%% -------------------------------------------------------------------------
\subsection{State Estimation}
\label{sec:state_estimation}
%% -------------------------------------------------------------------------

\subsubsection{IAP-RQ-010: 状态扩展：clk\_bias δt [m], clk\_drift δṫ [m/s] 加入 factor graph (`C(i)=gtsam::Vector2`)}
\label{ssec:iap_rq_010}

\textbf{Implementation:} \texttt{`include/iap/odometry/estimation\_frame.hpp`, `odometry\_estimation\_imu.hpp/.cpp`}.

\textbf{Verification:} `colcon build` 通过；`trace` 日志含 `clk\_bias/clk\_drift`.
\textbf{Logged metrics:} \texttt{`clk\_bias(m), clk\_drift(m/s)`}.
\textbf{Status:} **DONE**.

\subsubsection{IAP-RQ-015: Expose Σ\_p 位置协方差块（3×3）到 EstimationFrame；供 PL proxy 使用}
\label{ssec:iap_rq_015}

\textbf{Implementation:} \texttt{`estimation\_frame.hpp` (+sigma\_p), `odometry\_estimation\_imu.cpp` (marginalCovariance)}.

The position covariance proxy $\Sigma_p \in \mathbb{R}^{3\times3}$ is extracted
  from the fixed-lag smoother marginal:
  \begin{equation}
    \sigma_p = \sqrt{\lambda_{\max}(\Sigma_p)}, \qquad
    \mathrm{PL}_{\mathrm{proxy}} = K \sigma_p
    \label{eq:sigma_p}
  \end{equation}
  where $K=3$ corresponds to the $3\sigma$ coverage factor.

\textbf{Verification:} `trace` 日志含 `trace(Σ\_p)` 和 `lambda\_max(Σ\_p)`.
\textbf{Logged metrics:} \texttt{`trace\_sigma\_p, lambda\_max\_sigma\_p, PL\_proxy`}.
\textbf{Status:} **DONE**.

\subsubsection{IAP-RQ-020: GNSS 紧耦合观测：伪距 + 多普勒建因子（含 clock bias/drift，per-sat）}
\label{ssec:iap_rq_020}

\textbf{Implementation:} \texttt{`include/iap/gnss/`, `src/iap/gnss/` (PseudorangeFactor, DopplerFactor, GnssHandler)}.

The pseudorange residual for satellite $k$ is:
  \begin{equation}
    r_k^{\rho} = z_k^{\rho} - \bigl(\|\mathbf{p}_r - \mathbf{p}_k\| + \delta t\bigr),
    \label{eq:pseudorange}
  \end{equation}
  and the Doppler residual is:
  \begin{equation}
    r_k^{d} = z_k^{d} - \bigl(\mathbf{e}_k^\top(\mathbf{v}_r - \mathbf{v}_k) + \dot{\delta t}\bigr),
    \label{eq:doppler}
  \end{equation}
  where $\mathbf{e}_k = (\mathbf{p}_r - \mathbf{p}_k)/\|\cdot\|$ is the unit line-of-sight vector.

\textbf{Verification:} 关闭/开启 GNSS 对比.
\textbf{Logged metrics:} \texttt{`res\_pr, res\_dop, clk, clk\_dot`}.
\textbf{Status:} **DONE**.

\subsubsection{IAP-RQ-040: LiDAR ICP 因子健康度：退化/错配检测 → noise inflation / drop factor}
\label{ssec:iap_rq_040}

\textbf{Implementation:} \texttt{`EstimationFrame::IcpQuality`, `odometry\_estimation\_cpu.cpp` (Hessian cond, gamma\_lidar)}.

The LiDAR ICP condition number is computed from the Hessian block:
  \begin{equation}
    \kappa = \frac{\sigma_{\max}(H)}{\sigma_{\min}(H)}, \qquad
    \gamma_{\mathrm{lidar}} = \min\!\left(\sqrt{\kappa / \kappa_{\mathrm{th}}},\, \gamma_{\max}\right),
    \label{eq:icp_cond}
  \end{equation}
  and the LiDAR noise model is inflated by $\gamma_{\mathrm{lidar}}^{2}$ when $\kappa > \kappa_{\mathrm{th}}$.

\textbf{Verification:} 走廊/单面结构场景退化.
\textbf{Logged metrics:} \texttt{`icp\_rmse, inliers, cond, gamma\_lidar, drop`}.
\textbf{Status:} **DONE**.

\subsubsection{IAP-RQ-050: IMU 健康度：饱和/模型失配 → noise inflation / alarm}
\label{ssec:iap_rq_050}

\textbf{Implementation:} \texttt{`src/health/imu\_health.*`}.

\textbf{Verification:} 人为制造饱和或异常噪声.
\textbf{Logged metrics:} \texttt{`acc\_norm, gyro\_norm, sat\_flag, gamma\_imu`}.
\textbf{Status:} TODO.

%% -------------------------------------------------------------------------
\subsection{Environmental Perception}
\label{sec:environmental_perception}
%% -------------------------------------------------------------------------

\subsubsection{IAP-RQ-100: Trunk 检测与圆拟合（中心, 半径, confidence）}
\label{ssec:iap_rq_100}

\textbf{Implementation:} \texttt{`include/iap/trunk/trunk\_types.hpp`, `trunk\_detector.hpp/.cpp` (Kasa fit, grid BFS)}.

Trunk cylinders are fit using the Kasa circle estimator on a 2-D grid cluster.
  The TDOP for the detected trunk constellation is:
  \begin{equation}
    \mathrm{TDOP} = \sqrt{\mathrm{tr}\bigl((G^\top G)^{-1}\bigr)},
    \quad G_{k,:} = \mathbf{e}_{k,xy}^\top,
    \label{eq:tdop}
  \end{equation}
  where $\mathbf{e}_{k,xy}$ is the bearing to trunk $k$ in the horizontal plane.

\textbf{Verification:} 对进树林场景启动，日志输出 trunk 数量/半径/置信度.
\textbf{Logged metrics:} \texttt{`trunk\_count, radii, confidence[]`}.
\textbf{Status:} **DONE**.

\subsubsection{IAP-RQ-110: Trunk 健康度因子接口 (Baseline-A: 不入图)}
\label{ssec:iap_rq_110}

\textbf{Implementation:} \texttt{`TrunkDetector::health\_factor()` ([0,1])}.

\textbf{Verification:} health\textasciitilde{}0 时硬件应填充更大噪声.
\textbf{Logged metrics:} \texttt{`trunk\_health`}.
\textbf{Status:} **DONE (Baseline-A)**.

\subsubsection{IAP-RQ-120: TDOP 指标（角度多样性）}
\label{ssec:iap_rq_120}

\textbf{Implementation:} \texttt{`TrunkDetectionResult::tdop/tdop2/lambda\_min\_H`}.

\textbf{Verification:} 树更分散时 TDOP 下降.
\textbf{Logged metrics:} \texttt{`tdop, lambda\_min\_H`}.
\textbf{Status:} **DONE**.

%% -------------------------------------------------------------------------
\subsection{Integrity Monitoring}
\label{sec:integrity_monitoring}
%% -------------------------------------------------------------------------

\subsubsection{IAP-RQ-200: Integrity 输出：PL/AL/IM/mode + 关键中间量}
\label{ssec:iap_rq_200}

\textbf{Implementation:} \texttt{`include/iap/integrity/integrity\_types.hpp`, `integrity\_monitor.hpp/.cpp`}.

The protection level and alert limit are:
  \begin{equation}
    \mathrm{PL} = K_{\mathrm{pl}} \sqrt{\lambda_{\max}(\Sigma_p)}, \qquad
    \mathrm{AL} = \max\!\bigl(a_{\min},\, a_s \cdot d_{\mathrm{obs}} - r_{\mathrm{UAV}}\bigr),
    \label{eq:PL_AL}
  \end{equation}
  and the integrity margin is $\mathrm{IM} = \mathrm{AL} - \mathrm{PL}$.  The system is
  \emph{safe} when $\mathrm{IM} > 0$.

\textbf{Verification:} PL/AL/IM 曲线可画；mode 切换可复现.
\textbf{Logged metrics:} \texttt{`PL, AL, IM, mode`}.
\textbf{Status:} **DONE**.

\subsubsection{IAP-RQ-210: Alert Limit AL 由障碍距离动态给出}
\label{ssec:iap_rq_210}

\textbf{Implementation:} \texttt{`IntegrityMonitor::compute\_AL()`, `set\_obstacle\_distance()`}.

\textbf{Verification:} 越靠近障碍 AL 越小；日志可见.
\textbf{Logged metrics:} \texttt{`AL, obstacle\_dist`}.
\textbf{Status:} **DONE**.

\subsubsection{IAP-RQ-220: GNSS per-satellite NIS gating（RAIM-ish）}
\label{ssec:iap_rq_220}

\textbf{Implementation:} \texttt{`IntegrityMonitor::run\_gnss\_gating()` (chi2 test, gamma\_R, FDE greedy)}.

Per-satellite NIS gating (RAIM-ish):
  \begin{equation}
    \mathrm{NIS}_k = r_k^2 / \sigma_k^2 \overset{H_0}{\sim} \chi^2(1).
    \label{eq:nis}
  \end{equation}
  Satellite $k$ is excluded when $\mathrm{NIS}_k > \chi^2_{1,\alpha}$ with $\alpha=0.01$.

\textbf{Verification:} 注入 bias 卫星被降权/剔除.
\textbf{Logged metrics:} \texttt{`sat\_nis, gamma\_R, excluded\_sats`}.
\textbf{Status:} **DONE**.

%% -------------------------------------------------------------------------
\subsection{Trajectory Prediction}
\label{sec:trajectory_prediction}
%% -------------------------------------------------------------------------

\subsubsection{IAP-RQ-300: 候选轨迹生成（motion primitives）}
\label{ssec:iap_rq_300}

\textbf{Implementation:} \texttt{`include/iap/planner/trajectory\_types.hpp`, `trajectory\_generator.hpp/.cpp`}.

\textbf{Verification:} 生成 M 条候选轨迹并可视化时间戳点序列.
\textbf{Logged metrics:} \texttt{`trajectory\_count, speeds, yaw\_rates`}.
\textbf{Status:} **DONE**.

\subsubsection{IAP-RQ-310: 预测可见/可观测性集合（占位）}
\label{ssec:iap_rq_310}

\textbf{Implementation:} \texttt{`include/iap/planner/predicted\_integrity.hpp/.cpp` (placeholder)}.

\textbf{Verification:} TODO: 接地图后做 ray-check；当前返回占位值.
\textbf{Logged metrics:} \texttt{`n\_vis\_placeholder`}.
\textbf{Status:} **DONE (placeholder)**.

\subsubsection{IAP-RQ-320: 协方差传播 → Σ\_pred → PL\_pred}
\label{ssec:iap_rq_320}

\textbf{Implementation:} \texttt{`include/iap/planner/predicted\_integrity.hpp/.cpp` (sigma growth, K\_pl=3.0)}.

The isotropic covariance growth model predicts:
  \begin{equation}
    \sigma_{\mathrm{pred}}(t+\Delta t) = \sqrt{\sigma_{\mathrm{pred}}(t)^2 + \sigma_g^2 \Delta t},
    \qquad
    \mathrm{PL}_{\mathrm{pred}}(t) = K_{\mathrm{pl}} \cdot \sigma_{\mathrm{pred}}(t).
    \label{eq:sigma_grow}
  \end{equation}

\textbf{Verification:} PL\_pred 随时间增长且不同轨迹有差异；σ\_grow 可配置.
\textbf{Logged metrics:} \texttt{`PL\_pred(s), sigma\_pred(s)`}.
\textbf{Status:} **DONE (baseline)**.

%% -------------------------------------------------------------------------
\subsection{Integrity-Aware Planning}
\label{sec:integrity-aware_planning}
%% -------------------------------------------------------------------------

\subsubsection{IAP-RQ-400: Integrity-aware planning objective}
\label{ssec:iap_rq_400}

\textbf{Implementation:} \texttt{`include/iap/planner/integrity\_planner.hpp/.cpp`; `evaluate()`}.

The integrity-aware planning cost is:
  \begin{equation}
    J(\tau) = w_I \sum_{k} \bigl[\max(0,\, \mathrm{PL}_{\mathrm{pred},k} - \mathrm{AL})\bigr]^2
              + w_m \|\mathbf{p}_{N} - \mathbf{p}_{\mathrm{goal}}\|
              + w_s \sum_{k} \|\Delta \mathbf{v}_k\|,
    \label{eq:cost}
  \end{equation}
  where $w_I$ is boosted by a factor of $5$ in \textsc{search} mode.

\textbf{Verification:} IM<0时选绕行轨迹；J\_integrity > J\_goal场景可复现.
\textbf{Logged metrics:} \texttt{`J\_total, J\_integrity, J\_goal, J\_effort`}.
\textbf{Status:} **DONE**.

\subsubsection{IAP-RQ-410: Receding horizon loop}
\label{ssec:iap_rq_410}

\textbf{Implementation:} \texttt{`IntegrityPlanner::execution\_target()`, `plan()`}.

\textbf{Verification:} 调用`plan()`+`execution\_target()`模拟多步闭环.
\textbf{Logged metrics:} \texttt{`chosen\_traj\_id, dt\_execute`}.
\textbf{Status:} **DONE**.

%% -------------------------------------------------------------------------
\subsection{Experiments \& Metrics}
\label{sec:experiments_and_metrics}
%% -------------------------------------------------------------------------

\subsubsection{IAP-RQ-500: 三种 baseline（Passive/CovMin/IntegAware）}
\label{ssec:iap_rq_500}

\textbf{Implementation:} \texttt{`apps/iap\_experiment.cpp` (run\_baseline ×3)}.

\textbf{Verification:} 同场景三 baseline 均输出指标 CSV.
\textbf{Logged metrics:} \texttt{`baseline, violation\_frac, avg\_PL, mission\_success`}.
\textbf{Status:} **DONE (stub)**.

\subsubsection{IAP-RQ-510: 指标：Time(PL>AL)\%, AvgPL, MinIM, path/time/effort}
\label{ssec:iap_rq_510}

\textbf{Implementation:} \texttt{`include/iap/experiments/metrics.hpp` (MetricsCollector, write\_comparison\_table)}.

\textbf{Verification:} `ros2 run iap iap\_experiment` 输出 /tmp/*\_summary.md.
\textbf{Logged metrics:} \texttt{`violation\%, avg\_PL, min\_IM, path\_len, time, effort`}.
\textbf{Status:} **DONE (stub)**.

%% -------------------------------------------------------------------------
\subsection{Traceability Summary}
\label{sec:traceability}
%% -------------------------------------------------------------------------

Table~\ref{tab:trace} summarises the requirement-to-implementation mapping.

\begin{table*}[ht]
\caption{Requirement traceability matrix (auto-generated from \texttt{docs/TRACEABILITY.md})}
\label{tab:trace}
\centering
\small
\begin{tabular}{@{}lp{4.5cm}p{5.5cm}p{2.5cm}l@{}}
\toprule
Req ID & Description & Implementation & Metrics & Status \\
\midrule
\texttt{IAP-RQ-000} & Repo guardrails：AGENTS.md、doc-guard（pre-commit hook + tools/doc\_guard.py）、docs 三件套 & \texttt{`AGENTS.md`, `.githooks/pre-commit`, `tools/doc\_guard.py`, `docs/`} & \texttt{hook exit code} & **DONE** \\
\texttt{IAP-RQ-001} & Rename ROS2 package to `iap` & \texttt{`package.xml`, `CMakeLists.txt`, `src/iap/`, `include/iap/`, `cmake/iap-config.cmake.in`} & \texttt{grep iap` 可见} & build exit code 0 \\
\texttt{IAP-RQ-002} & Build artifacts: compile\_commands.json + 最小 demo & \texttt{`apps/iap\_status.cpp`, `launch/iap\_demo.launch.py`, `.clangd`, `CMakeLists.txt` (IAP\_VERSION, install launch/)} & \texttt{`iap\_status: OK`} & **DONE** \\
\texttt{IAP-RQ-010} & 状态扩展：clk\_bias δt [m], clk\_drift δṫ [m/s] 加入 factor graph (`C(i)=gtsam::Vector2`) & \texttt{`include/iap/odometry/estimation\_frame.hpp`, `odometry\_estimation\_imu.hpp/.cpp`} & \texttt{`clk\_bias(m), clk\_drift(m/s)`} & **DONE** \\
\texttt{IAP-RQ-015} & Expose Σ\_p 位置协方差块（3×3）到 EstimationFrame；供 PL proxy 使用 & \texttt{`estimation\_frame.hpp` (+sigma\_p), `odometry\_estimation\_imu.cpp` (marginalCovariance)} & \texttt{`trace\_sigma\_p, lambda\_max\_sigma\_p, PL\_proxy`} & **DONE** \\
\texttt{IAP-RQ-020} & GNSS 紧耦合观测：伪距 + 多普勒建因子（含 clock bias/drift，per-sat） & \texttt{`include/iap/gnss/`, `src/iap/gnss/` (PseudorangeFactor, DopplerFactor, GnssHandler)} & \texttt{`res\_pr, res\_dop, clk, clk\_dot`} & **DONE** \\
\texttt{IAP-RQ-030} & GNSS per-satellite NIS gating：downweight / exclude（RAIM-ish baseline） & \texttt{`src/integrity/gnss\_integrity.*`} & \texttt{`sat\_nis, exclude\_sats, gamma\_R, global\_nis`} & TODO \\
\texttt{IAP-RQ-040} & LiDAR ICP 因子健康度：退化/错配检测 → noise inflation / drop factor & \texttt{`EstimationFrame::IcpQuality`, `odometry\_estimation\_cpu.cpp` (Hessian cond, gamma\_lidar)} & \texttt{`icp\_rmse, inliers, cond, gamma\_lidar, drop`} & **DONE** \\
\texttt{IAP-RQ-100} & Trunk 检测与圆拟合（中心, 半径, confidence） & \texttt{`include/iap/trunk/trunk\_types.hpp`, `trunk\_detector.hpp/.cpp` (Kasa fit, grid BFS)} & \texttt{`trunk\_count, radii, confidence[]`} & **DONE** \\
\texttt{IAP-RQ-110} & Trunk 健康度因子接口 (Baseline-A: 不入图) & \texttt{`TrunkDetector::health\_factor()` ([0,1])} & \texttt{`trunk\_health`} & **DONE (Baseline-A)** \\
\texttt{IAP-RQ-120} & TDOP 指标（角度多样性） & \texttt{`TrunkDetectionResult::tdop/tdop2/lambda\_min\_H`} & \texttt{`tdop, lambda\_min\_H`} & **DONE** \\
\texttt{IAP-RQ-200} & Integrity 输出：PL/AL/IM/mode + 关键中间量 & \texttt{`include/iap/integrity/integrity\_types.hpp`, `integrity\_monitor.hpp/.cpp`} & \texttt{`PL, AL, IM, mode`} & **DONE** \\
\texttt{IAP-RQ-210} & Alert Limit AL 由障碍距离动态给出 & \texttt{`IntegrityMonitor::compute\_AL()`, `set\_obstacle\_distance()`} & \texttt{`AL, obstacle\_dist`} & **DONE** \\
\texttt{IAP-RQ-220} & GNSS per-satellite NIS gating（RAIM-ish） & \texttt{`IntegrityMonitor::run\_gnss\_gating()` (chi2 test, gamma\_R, FDE greedy)} & \texttt{`sat\_nis, gamma\_R, excluded\_sats`} & **DONE** \\
\texttt{IAP-RQ-300} & 候选轨迹生成（motion primitives） & \texttt{`include/iap/planner/trajectory\_types.hpp`, `trajectory\_generator.hpp/.cpp`} & \texttt{`trajectory\_count, speeds, yaw\_rates`} & **DONE** \\
\texttt{IAP-RQ-310} & 预测可见/可观测性集合（占位） & \texttt{`include/iap/planner/predicted\_integrity.hpp/.cpp` (placeholder)} & \texttt{`n\_vis\_placeholder`} & **DONE (placeholder)** \\
\texttt{IAP-RQ-320} & 协方差传播 → Σ\_pred → PL\_pred & \texttt{`include/iap/planner/predicted\_integrity.hpp/.cpp` (sigma growth, K\_pl=3.0)} & \texttt{`PL\_pred(s), sigma\_pred(s)`} & **DONE (baseline)** \\
\texttt{IAP-RQ-400} & Integrity-aware planning objective & \texttt{`include/iap/planner/integrity\_planner.hpp/.cpp`; `evaluate()`} & \texttt{`J\_total, J\_integrity, J\_goal, J\_effort`} & **DONE** \\
\texttt{IAP-RQ-410} & Receding horizon loop & \texttt{`IntegrityPlanner::execution\_target()`, `plan()`} & \texttt{`chosen\_traj\_id, dt\_execute`} & **DONE** \\
\texttt{IAP-RQ-500} & 三种 baseline（Passive/CovMin/IntegAware） & \texttt{`apps/iap\_experiment.cpp` (run\_baseline ×3)} & \texttt{`baseline, violation\_frac, avg\_PL, mission\_success`} & **DONE (stub)** \\
\texttt{IAP-RQ-510} & 指标：Time(PL>AL)\%, AvgPL, MinIM, path/time/effort & \texttt{`include/iap/experiments/metrics.hpp` (MetricsCollector, write\_comparison\_table)} & \texttt{`violation\%, avg\_PL, min\_IM, path\_len, time, effort`} & **DONE (stub)** \\
\texttt{IAP-RQ-050} & IMU 健康度：饱和/模型失配 → noise inflation / alarm & \texttt{`src/health/imu\_health.*`} & \texttt{`acc\_norm, gyro\_norm, sat\_flag, gamma\_imu`} & TODO \\
\texttt{IAP-RQ-060} & 规划目标：J(τ)=Σ hinge(PL\_pred-AL)\textasciicircum{}2 + λ\_goal*dist + λ\_u*effort & \texttt{`src/planner/cost.*`} & \texttt{`J\_total, J\_integrity, J\_goal, J\_effort`} & TODO \\
\texttt{IAP-RQ-070} & 预测层：对候选轨迹预测 PL\_pred（baseline：代理；升级：ARAIM） & \texttt{`src/predictor/*`} & \texttt{`PL\_pred(s), AL(s), IM\_pred(s)`} & TODO \\
\texttt{IAP-RQ-080} & Receding horizon 闭环：执行第一段，重估计重规划 & \texttt{`src/planner/mpc\_loop.*` / `apps/*`} & \texttt{`chosen\_traj\_id, replanning\_rate`} & TODO \\
\texttt{IAP-RQ-090} & 实验与消融：直飞 vs 协方差最小 vs 完整性驱动 & \texttt{`apps/experiments/*` + `docs/*`} & \texttt{`violation\_time, min\_IM, path\_len, time`} & TODO \\
\bottomrule
\end{tabular}
\end{table*}

\end{document}
